\subsection{Трассировка лучей методом Монте-Карло}

Если учитывать рассеивание среды, то явно получить решение не представляется возможным. Поэтому метод Монте-Карло честно просчитывает поведение луча.

Простейшая идея заключается в том, чтобы испускать из всех источников света лучи, смотреть как они будут себя вести, и если они попадают на сенсор наблюдения, накапливать их, и таким образом получить изображение. Однако, вероятность случайно наблюдаемого луча попасть на сенсор крайне мала. Поэтому куда практичней идти от конца луча к его началу, в этом и заключается метод Монте-Карло.

В уравнении \ref{eq:graphics_rte} удобней выделить потери и прирост света:
%
\[
	\frac{\mathrm dL(\mathrm p, \omega)}{\mathrm dt} = 
	-\sigma_t(\mathrm p, \omega) L(\mathrm p, \omega) +
	S(\mathrm p, \omega),
\]
%
где
%
\[
	S(\mathrm p, \omega) =
	\sigma_a(\mathrm p, \omega) L_e(\mathrm p, \omega) + 
	\sigma_s(\mathrm p, \omega) \int_{S^2} p(\mathrm p, \omega', \omega) L(\mathrm p, \omega') \mathrm d\omega'.
\]
%
По аналогии с непрозрачностью, введенной для получения уравнения \ref{eq:int_rte_emm_abs}, вводится коэффициент пропускания
%
\[
	T(\mathrm p_1, \mathrm p_2) =
	e^{-\int_{\mathrm p_1}^{\mathrm p_2} \sigma_t(s) \mathrm ds}.
\]

Рассматривается конкретный луч $\mathrm r(t) = \mathrm o + t\omega$). Для того, чтобы не нагромождать формулы, далее $L(t, \omega)=L(r(t), \omega)$. Итоговое интегральное решение (доказательство приводится в \cite{max1995optical}) имеет вид:
%
\begin{align*}
	L(0, \omega) = T(0, d) L(d, \omega) + \int_0^d T(0, t) S(t, \omega) \mathrm dt, \\
	S(t, \omega) =
	\sigma_a(t, \omega) L_e(t, \omega) +
	\sigma_s(t, \omega) \int_{S^2} \mathrm p(t, \omega', \omega) L(t, \omega') \mathrm d\omega'.
\end{align*}
%
Удобно за $d$ принять точку, яркость которой известна. Простейший случай, когда заранее задано освещение от фона, например $f$. Это просто некоторая функция на сфере, и тогда в достаточно далекой точке $d=\infty$ вне объема $L(\infty, \omega)=f(\omega)$. Это позволяет считать слагаемое $T(0, d) L(d, \omega)$ известным (или по крайней мере легко считаемым). Также можно посчитать и слагаемое
%
\[
	\int_0^\infty T(0, t) \sigma_a(t, \omega) L_e(t, \omega) \mathrm dt,
\]
%
применив как и в случае с ray casting конечные суммы. Метод Монте-Карло предлагает способ вычисления последнего, самого сложного слагаемого
%
\begin{equation}
	\label{eq:scattering_integral}
	\int_0^d T(0, t) \sigma_s(t, \omega) \mathrm dt
	\int_{S^2} \mathrm p(t, \omega', \omega) \mathrm d\omega',
\end{equation}
%
и заключается в приближении интеграла
%
\[
	\int_a^b f(x) \mathrm dx \approx
	\frac1n \sum_{i=1}^n \frac{f(X_i)}{P(X_i)},
\]
%
где $X_i$ сэмплированы из распределения с плотностью $P$\footnote{Большая буква для плотности распределения выбрана, чтобы не путать его с фазовой функцией.}.

При вычислении интеграла \ref{eq:scattering_integral} удобно сэмплировать точки внутри объёма из распределения с плотностью $P(t) = \sigma_t(t) T(0, t)$, за $n$ сэмплов из которого получится приближение
%
\[
	\sum_{i=1}^n \frac{\sigma_s (t_i, \omega)}{\sigma_t (t_i, \omega)}
	\int_{S^2} \mathrm p(t_i, \omega', \omega) L(t_i, \omega') \mathrm d\omega'.
\]
%
Второй интеграл можно приблизить аналогичным образом, только сэмплировать всего 1\footnote{Можно сэмплировать и больше, но на практике так не делают} перенаправленный луч по направлению $\omega$ из распределения, задаваемого фазовой функцией, тогда приближение примет итоговый вид
%
\[
	\sum_{i=1}^n \frac{\sigma_s (t_i, \omega)}{\sigma_t (t_i, \omega)} L(t_i, \omega'),
\]
%
где $t_i \sim P(t_i)$, $\omega' \sim p(t_i, \omega', \omega)$.

Как видно, это требует уже вычисления $L(t_i, \omega')$ аналогичным способом, который как раз и запускает рекурсию. В зависимости от того, насколько глубоко разрешена рекурсия --- будут получаться рендеры различной физической корректности.